\tikzset{>=latex} % for LaTeX arrow head
\contourlength{0.8pt}

\colorlet{xcol}{blue!70!black}
\colorlet{myred}{red!65!black}
\tikzstyle{rvec}=[->,xcol,very thick,line cap=round]
\tikzstyle{mass line}=[line width=0.5,draw=red!30!black]
\tikzstyle{myarr}=[-{Latex[length=3,width=2]},blue!40!black]
\tikzstyle{myarr2}=[{Latex[length=3,width=2]}-{Latex[length=3,width=2]},blue!40!black]
\tikzstyle{mass}=[mass line, %rounded corners=1,
                  top color=red!40!black!30,bottom color=red!40!black!10,shading angle=30]
\tikzstyle{middle mass}=[mass line,top color=red!40!black!50,bottom color=red!40!black!50,
                         middle color=red!40!black!10,shading angle=30]

\def\r{0.05} % pulley small radius
\tikzset{
  pics/rotarr/.style={
    code={
      \draw[white,line width=0.8] ({#1*cos(210)},0) arc(-210:35:{#1} and {0.35*#1});
      \draw[-{>[flex'=1]}] ({#1*cos(210)},0) coordinate (W1) arc(-210:35:{#1} and {0.35*#1})
        node[midway] (W2) {} --++ (150:0.1) coordinate (W3);
  }},
  pics/rotarr/.default=0.3,
}

\LARGE

% MOMENT OF INERTIA - masses on rods
\def\ang{-20} % angle of whole picture




% MOMENT OF INERTIA - ROD (axis through end)


\begin{tikzpicture}[rotate=\ang]
  \rod
  \draw[line cap=round] (\L/2,-0.8*\h) --++ (0,2*\h);
  \pic[xscale=1,rotate=\ang] at (\L/2,0.7*\h) {rotarr};
  \node[right=0] at (W3) {$\omega$};
  \draw[white,line width=0.7] (BL)++(0,-1.8*\Ry) --++ (\L,0);
  \draw[myarr2] (BL)++(0,-1.8*\Ry) --++ (\L,0)
    node[midway,below=-1,scale=0.9] {$L$};
\end{tikzpicture}
